\section{Learning Objectives}

After completing this chapter, readers will be able to:
\begin{itemize}
    \item Design and implement Terraform modules for ML infrastructure
    \item Deploy ML workloads on Kubernetes using Helm charts
    \item Manage cloud resources across AWS, GCP, and Azure
    \item Orchestrate containerized ML workloads at scale
    \item Implement disaster recovery and backup strategies for ML systems
    \item Optimize infrastructure costs through Infrastructure as Code
    \item Automate infrastructure provisioning and management
    \item Apply best practices for production ML infrastructure
\end{itemize}

\section{Introduction}

Infrastructure as Code (IaC) is the practice of managing and provisioning infrastructure through machine-readable definition files rather than manual processes. For ML systems, IaC provides reproducibility, version control, and automation of complex infrastructure requirements.

\textbf{Why IaC for Machine Learning:}

\begin{itemize}
    \item \textbf{Reproducibility:} Recreate exact infrastructure configurations across environments
    \item \textbf{Version Control:} Track infrastructure changes alongside code changes
    \item \textbf{Automation:} Eliminate manual provisioning and reduce human error
    \item \textbf{Scalability:} Rapidly scale infrastructure up or down based on demand
    \item \textbf{Cost Optimization:} Programmatically manage resources to minimize costs
    \item \textbf{Disaster Recovery:} Quickly rebuild infrastructure from code
    \item \textbf{Collaboration:} Share infrastructure definitions across teams
    \item \textbf{Testing:} Test infrastructure changes before production deployment
\end{itemize}

\textbf{ML Infrastructure Components:}

\begin{enumerate}
    \item \textbf{Compute:} GPU/CPU instances for training and inference
    \item \textbf{Storage:} Data lakes, model registries, artifact storage
    \item \textbf{Networking:} VPCs, load balancers, API gateways
    \item \textbf{Orchestration:} Kubernetes clusters, container registries
    \item \textbf{Data:} Databases, data warehouses, streaming platforms
    \item \textbf{Monitoring:} Logging, metrics, alerting infrastructure
    \item \textbf{Security:} IAM, secrets management, encryption
\end{enumerate}

\section{Terraform Modules for ML Infrastructure}

Terraform is an open-source IaC tool that enables declarative infrastructure management across multiple cloud providers.

\subsection{Core ML Infrastructure Module}

\begin{lstlisting}[language=Python, caption=Terraform Module Structure for ML Platform]
# modules/ml-platform/main.tf
"""
Root module for ML platform infrastructure
Orchestrates compute, storage, networking, and monitoring
"""

terraform {
  required_version = ">= 1.5.0"

  required_providers {
    aws = {
      source  = "hashicorp/aws"
      version = "~> 5.0"
    }
    kubernetes = {
      source  = "hashicorp/kubernetes"
      version = "~> 2.23"
    }
  }
}

# Local variables
locals {
  common_tags = {
    Project     = var.project_name
    Environment = var.environment
    ManagedBy   = "Terraform"
    CostCenter  = var.cost_center
    Owner       = var.owner_email
  }

  name_prefix = "${var.project_name}-${var.environment}"
}

# Data sources
data "aws_caller_identity" "current" {}
data "aws_region" "current" {}
data "aws_availability_zones" "available" {
  state = "available"
}

# VPC Module
module "vpc" {
  source  = "terraform-aws-modules/vpc/aws"
  version = "~> 5.0"

  name = "${local.name_prefix}-vpc"
  cidr = var.vpc_cidr

  azs             = slice(data.aws_availability_zones.available.names, 0, 3)
  private_subnets = [for i in range(3) : cidrsubnet(var.vpc_cidr, 8, i)]
  public_subnets  = [for i in range(3) : cidrsubnet(var.vpc_cidr, 8, i + 100)]

  enable_nat_gateway   = true
  single_nat_gateway   = var.environment != "production"
  enable_dns_hostnames = true
  enable_dns_support   = true

  # Enable VPC Flow Logs
  enable_flow_log                      = true
  create_flow_log_cloudwatch_iam_role  = true
  create_flow_log_cloudwatch_log_group = true

  tags = local.common_tags
}

# EKS Cluster Module
module "eks" {
  source  = "terraform-aws-modules/eks/aws"
  version = "~> 19.0"

  cluster_name    = "${local.name_prefix}-eks"
  cluster_version = var.kubernetes_version

  vpc_id                   = module.vpc.vpc_id
  subnet_ids               = module.vpc.private_subnets
  control_plane_subnet_ids = module.vpc.private_subnets

  # Cluster endpoint access
  cluster_endpoint_public_access  = var.environment != "production"
  cluster_endpoint_private_access = true

  # Cluster add-ons
  cluster_addons = {
    coredns = {
      most_recent = true
    }
    kube-proxy = {
      most_recent = true
    }
    vpc-cni = {
      most_recent = true
    }
    aws-ebs-csi-driver = {
      most_recent = true
    }
  }

  # Managed node groups
  eks_managed_node_groups = {
    # CPU workloads
    general = {
      name = "${local.name_prefix}-general"

      instance_types = ["m5.xlarge"]
      capacity_type  = "ON_DEMAND"

      min_size     = var.environment == "production" ? 2 : 1
      max_size     = 10
      desired_size = var.environment == "production" ? 3 : 2

      labels = {
        workload = "general"
      }

      taints = []

      tags = merge(local.common_tags, {
        NodeGroup = "general"
      })
    }

    # GPU workloads for training
    gpu_training = {
      name = "${local.name_prefix}-gpu-training"

      instance_types = ["p3.2xlarge"]
      capacity_type  = var.use_spot_instances ? "SPOT" : "ON_DEMAND"

      min_size     = 0
      max_size     = 20
      desired_size = 0  # Scale from zero

      labels = {
        workload    = "training"
        gpu         = "nvidia-v100"
        accelerator = "true"
      }

      taints = [{
        key    = "nvidia.com/gpu"
        value  = "true"
        effect = "NO_SCHEDULE"
      }]

      # Install NVIDIA device plugin
      pre_bootstrap_user_data = <<-EOT
        #!/bin/bash
        /etc/eks/bootstrap.sh ${local.name_prefix}-eks \
          --kubelet-extra-args '--node-labels=workload=training,gpu=nvidia-v100'
      EOT

      tags = merge(local.common_tags, {
        NodeGroup = "gpu-training"
      })
    }

    # Inference workloads
    inference = {
      name = "${local.name_prefix}-inference"

      instance_types = ["c5.2xlarge"]
      capacity_type  = "ON_DEMAND"

      min_size     = var.environment == "production" ? 2 : 1
      max_size     = 50
      desired_size = var.environment == "production" ? 3 : 1

      labels = {
        workload = "inference"
      }

      tags = merge(local.common_tags, {
        NodeGroup = "inference"
      })
    }
  }

  # Cluster security group rules
  cluster_security_group_additional_rules = {
    ingress_nodes_ephemeral_ports_tcp = {
      description                = "Nodes on ephemeral ports"
      protocol                   = "tcp"
      from_port                  = 1025
      to_port                    = 65535
      type                       = "ingress"
      source_node_security_group = true
    }
  }

  # Node security group rules
  node_security_group_additional_rules = {
    ingress_self_all = {
      description = "Node to node all ports/protocols"
      protocol    = "-1"
      from_port   = 0
      to_port     = 0
      type        = "ingress"
      self        = true
    }
    egress_all = {
      description      = "Node all egress"
      protocol         = "-1"
      from_port        = 0
      to_port          = 0
      type             = "egress"
      cidr_blocks      = ["0.0.0.0/0"]
      ipv6_cidr_blocks = ["::/0"]
    }
  }

  tags = local.common_tags
}

# S3 Buckets for ML Assets
module "s3_buckets" {
  source = "./modules/s3-buckets"

  project_name = var.project_name
  environment  = var.environment

  buckets = {
    data = {
      name        = "${local.name_prefix}-data"
      versioning  = true
      encryption  = true
      lifecycle_rules = [
        {
          id      = "archive_old_data"
          enabled = true
          transitions = [
            {
              days          = 90
              storage_class = "INTELLIGENT_TIERING"
            },
            {
              days          = 180
              storage_class = "GLACIER"
            }
          ]
        }
      ]
    }

    models = {
      name        = "${local.name_prefix}-models"
      versioning  = true
      encryption  = true
      lifecycle_rules = [
        {
          id      = "keep_versions"
          enabled = true
          noncurrent_version_expiration = {
            days = 90
          }
        }
      ]
    }

    artifacts = {
      name        = "${local.name_prefix}-artifacts"
      versioning  = false
      encryption  = true
      lifecycle_rules = [
        {
          id      = "expire_artifacts"
          enabled = true
          expiration = {
            days = 30
          }
        }
      ]
    }

    logs = {
      name        = "${local.name_prefix}-logs"
      versioning  = false
      encryption  = true
      lifecycle_rules = [
        {
          id      = "expire_logs"
          enabled = true
          expiration = {
            days = 7
          }
        }
      ]
    }
  }

  tags = local.common_tags
}

# RDS for Metadata Storage
module "rds" {
  source = "./modules/rds"

  identifier = "${local.name_prefix}-metadata"

  engine               = "postgres"
  engine_version       = "15.3"
  instance_class       = var.db_instance_type
  allocated_storage    = 100
  max_allocated_storage = 1000
  storage_encrypted    = true

  db_name  = "mlmetadata"
  username = var.db_master_username

  vpc_id             = module.vpc.vpc_id
  subnet_ids         = module.vpc.private_subnets

  backup_retention_period = var.environment == "production" ? 30 : 7
  backup_window          = "03:00-04:00"
  maintenance_window     = "mon:04:00-mon:05:00"

  deletion_protection = var.environment == "production"
  skip_final_snapshot = var.environment != "production"

  # Performance Insights
  performance_insights_enabled    = true
  performance_insights_retention_period = 7

  # Enhanced monitoring
  enabled_cloudwatch_logs_exports = ["postgresql", "upgrade"]
  monitoring_interval            = 60

  tags = local.common_tags
}

# ElastiCache for Caching
module "elasticache" {
  source = "./modules/elasticache"

  cluster_id = "${local.name_prefix}-redis"

  engine         = "redis"
  engine_version = "7.0"
  node_type      = var.cache_node_type
  num_cache_nodes = var.environment == "production" ? 2 : 1

  vpc_id     = module.vpc.vpc_id
  subnet_ids = module.vpc.private_subnets

  automatic_failover_enabled = var.environment == "production"
  multi_az_enabled          = var.environment == "production"

  at_rest_encryption_enabled = true
  transit_encryption_enabled = true

  snapshot_retention_limit = var.environment == "production" ? 5 : 1
  snapshot_window         = "03:00-05:00"

  tags = local.common_tags
}

# ECR Repositories
resource "aws_ecr_repository" "ml_images" {
  for_each = toset([
    "training",
    "inference",
    "preprocessing",
    "monitoring"
  ])

  name                 = "${local.name_prefix}-${each.key}"
  image_tag_mutability = "MUTABLE"

  image_scanning_configuration {
    scan_on_push = true
  }

  encryption_configuration {
    encryption_type = "AES256"
  }

  tags = local.common_tags
}

resource "aws_ecr_lifecycle_policy" "ml_images" {
  for_each = aws_ecr_repository.ml_images

  repository = each.value.name

  policy = jsonencode({
    rules = [
      {
        rulePriority = 1
        description  = "Keep last 30 images"
        selection = {
          tagStatus     = "any"
          countType     = "imageCountMoreThan"
          countNumber   = 30
        }
        action = {
          type = "expire"
        }
      },
      {
        rulePriority = 2
        description  = "Remove untagged images after 7 days"
        selection = {
          tagStatus   = "untagged"
          countType   = "sinceImagePushed"
          countUnit   = "days"
          countNumber = 7
        }
        action = {
          type = "expire"
        }
      }
    ]
  })
}

# SageMaker Domain for Notebooks
resource "aws_sagemaker_domain" "ml_studio" {
  domain_name = "${local.name_prefix}-studio"
  auth_mode   = "IAM"
  vpc_id      = module.vpc.vpc_id
  subnet_ids  = module.vpc.private_subnets

  default_user_settings {
    execution_role = aws_iam_role.sagemaker_execution.arn

    jupyter_server_app_settings {
      default_resource_spec {
        instance_type       = "ml.t3.medium"
        sagemaker_image_arn = "arn:aws:sagemaker:${data.aws_region.current.name}:${data.aws_caller_identity.current.account_id}:image/jupyter-server-3"
      }
    }

    kernel_gateway_app_settings {
      default_resource_spec {
        instance_type       = "ml.t3.medium"
        sagemaker_image_arn = "arn:aws:sagemaker:${data.aws_region.current.name}:${data.aws_caller_identity.current.account_id}:image/datascience-1.0"
      }
    }
  }

  tags = local.common_tags
}

# IAM Role for SageMaker
resource "aws_iam_role" "sagemaker_execution" {
  name = "${local.name_prefix}-sagemaker-execution"

  assume_role_policy = jsonencode({
    Version = "2012-10-17"
    Statement = [
      {
        Action = "sts:AssumeRole"
        Effect = "Allow"
        Principal = {
          Service = "sagemaker.amazonaws.com"
        }
      }
    ]
  })

  tags = local.common_tags
}

resource "aws_iam_role_policy_attachment" "sagemaker_full_access" {
  role       = aws_iam_role.sagemaker_execution.name
  policy_arn = "arn:aws:iam::aws:policy/AmazonSageMakerFullAccess"
}

resource "aws_iam_role_policy" "sagemaker_s3_access" {
  name = "${local.name_prefix}-sagemaker-s3"
  role = aws_iam_role.sagemaker_execution.id

  policy = jsonencode({
    Version = "2012-10-17"
    Statement = [
      {
        Effect = "Allow"
        Action = [
          "s3:GetObject",
          "s3:PutObject",
          "s3:DeleteObject",
          "s3:ListBucket"
        ]
        Resource = [
          for bucket in module.s3_buckets.bucket_arns : "${bucket}/*"
        ]
      }
    ]
  })
}

# CloudWatch Log Groups
resource "aws_cloudwatch_log_group" "ml_logs" {
  for_each = toset([
    "training",
    "inference",
    "preprocessing",
    "monitoring"
  ])

  name              = "/aws/ml/${local.name_prefix}/${each.key}"
  retention_in_days = var.log_retention_days

  tags = local.common_tags
}

# KMS Key for Encryption
resource "aws_kms_key" "ml_data" {
  description             = "KMS key for ML data encryption"
  deletion_window_in_days = 30
  enable_key_rotation     = true

  tags = merge(local.common_tags, {
    Purpose = "ML Data Encryption"
  })
}

resource "aws_kms_alias" "ml_data" {
  name          = "alias/${local.name_prefix}-ml-data"
  target_key_id = aws_kms_key.ml_data.key_id
}
\end{lstlisting}

\subsection{Terraform Variables and Outputs}

\begin{lstlisting}[language=Python, caption=Module Variables Definition]
# modules/ml-platform/variables.tf
"""
Input variables for ML platform module
"""

variable "project_name" {
  description = "Name of the ML project"
  type        = string

  validation {
    condition     = can(regex("^[a-z0-9-]+$", var.project_name))
    error_message = "Project name must contain only lowercase letters, numbers, and hyphens."
  }
}

variable "environment" {
  description = "Environment name (development, staging, production)"
  type        = string

  validation {
    condition     = contains(["development", "staging", "production"], var.environment)
    error_message = "Environment must be development, staging, or production."
  }
}

variable "owner_email" {
  description = "Email of the infrastructure owner"
  type        = string

  validation {
    condition     = can(regex("^[a-zA-Z0-9._%+-]+@[a-zA-Z0-9.-]+\\.[a-zA-Z]{2,}$", var.owner_email))
    error_message = "Must be a valid email address."
  }
}

variable "cost_center" {
  description = "Cost center for billing"
  type        = string
  default     = "ml-platform"
}

variable "vpc_cidr" {
  description = "CIDR block for VPC"
  type        = string
  default     = "10.0.0.0/16"

  validation {
    condition     = can(cidrhost(var.vpc_cidr, 0))
    error_message = "Must be a valid CIDR block."
  }
}

variable "kubernetes_version" {
  description = "Kubernetes version for EKS cluster"
  type        = string
  default     = "1.28"
}

variable "use_spot_instances" {
  description = "Use spot instances for GPU nodes"
  type        = bool
  default     = false
}

variable "db_instance_type" {
  description = "RDS instance type"
  type        = string
  default     = "db.t3.large"
}

variable "db_master_username" {
  description = "Master username for RDS"
  type        = string
  default     = "admin"
  sensitive   = true
}

variable "cache_node_type" {
  description = "ElastiCache node type"
  type        = string
  default     = "cache.t3.medium"
}

variable "log_retention_days" {
  description = "CloudWatch log retention in days"
  type        = number
  default     = 7

  validation {
    condition     = contains([1, 3, 5, 7, 14, 30, 60, 90, 120, 150, 180, 365, 400, 545, 731, 1827, 3653], var.log_retention_days)
    error_message = "Must be a valid CloudWatch log retention period."
  }
}

variable "enable_monitoring" {
  description = "Enable enhanced monitoring"
  type        = bool
  default     = true
}

variable "enable_auto_scaling" {
  description = "Enable cluster autoscaling"
  type        = bool
  default     = true
}
\end{lstlisting}

\begin{lstlisting}[language=Python, caption=Module Outputs]
# modules/ml-platform/outputs.tf
"""
Outputs from ML platform module
"""

output "vpc_id" {
  description = "ID of the VPC"
  value       = module.vpc.vpc_id
}

output "vpc_cidr" {
  description = "CIDR block of the VPC"
  value       = module.vpc.vpc_cidr_block
}

output "private_subnets" {
  description = "IDs of private subnets"
  value       = module.vpc.private_subnets
}

output "public_subnets" {
  description = "IDs of public subnets"
  value       = module.vpc.public_subnets
}

output "eks_cluster_id" {
  description = "EKS cluster ID"
  value       = module.eks.cluster_id
}

output "eks_cluster_endpoint" {
  description = "Endpoint for EKS control plane"
  value       = module.eks.cluster_endpoint
}

output "eks_cluster_security_group_id" {
  description = "Security group ID attached to the EKS cluster"
  value       = module.eks.cluster_security_group_id
}

output "eks_cluster_certificate_authority_data" {
  description = "Base64 encoded certificate data required to communicate with the cluster"
  value       = module.eks.cluster_certificate_authority_data
  sensitive   = true
}

output "eks_oidc_provider_arn" {
  description = "ARN of the OIDC Provider for EKS"
  value       = module.eks.oidc_provider_arn
}

output "s3_bucket_names" {
  description = "Names of S3 buckets"
  value       = module.s3_buckets.bucket_names
}

output "s3_bucket_arns" {
  description = "ARNs of S3 buckets"
  value       = module.s3_buckets.bucket_arns
}

output "rds_endpoint" {
  description = "RDS instance endpoint"
  value       = module.rds.endpoint
  sensitive   = true
}

output "rds_database_name" {
  description = "Name of the RDS database"
  value       = module.rds.database_name
}

output "redis_endpoint" {
  description = "ElastiCache Redis endpoint"
  value       = module.elasticache.endpoint
  sensitive   = true
}

output "ecr_repository_urls" {
  description = "URLs of ECR repositories"
  value = {
    for k, v in aws_ecr_repository.ml_images : k => v.repository_url
  }
}

output "sagemaker_domain_id" {
  description = "SageMaker domain ID"
  value       = aws_sagemaker_domain.ml_studio.id
}

output "sagemaker_execution_role_arn" {
  description = "ARN of SageMaker execution role"
  value       = aws_iam_role.sagemaker_execution.arn
}

output "kms_key_id" {
  description = "ID of KMS key for encryption"
  value       = aws_kms_key.ml_data.key_id
}

output "kms_key_arn" {
  description = "ARN of KMS key for encryption"
  value       = aws_kms_key.ml_data.arn
}
\end{lstlisting}

\subsection{Environment-Specific Configuration}

\begin{lstlisting}[language=Python, caption=Production Environment Configuration]
# environments/production/main.tf
"""
Production environment configuration
"""

terraform {
  backend "s3" {
    bucket         = "ml-platform-terraform-state"
    key            = "production/terraform.tfstate"
    region         = "us-west-2"
    encrypt        = true
    dynamodb_table = "terraform-state-lock"
  }
}

provider "aws" {
  region = "us-west-2"

  default_tags {
    tags = {
      Environment = "production"
      ManagedBy   = "Terraform"
      Project     = "ml-platform"
    }
  }
}

module "ml_platform" {
  source = "../../modules/ml-platform"

  project_name = "ml-platform"
  environment  = "production"
  owner_email  = "ml-team@company.com"
  cost_center  = "ml-engineering"

  # Network configuration
  vpc_cidr = "10.0.0.0/16"

  # Kubernetes configuration
  kubernetes_version  = "1.28"
  use_spot_instances = false  # Use on-demand for production
  enable_auto_scaling = true

  # Database configuration
  db_instance_type    = "db.r6g.xlarge"
  db_master_username  = var.db_master_username

  # Cache configuration
  cache_node_type = "cache.r6g.large"

  # Monitoring and logging
  log_retention_days = 30
  enable_monitoring  = true
}

# Outputs
output "eks_cluster_endpoint" {
  value       = module.ml_platform.eks_cluster_endpoint
  description = "EKS cluster endpoint"
}

output "model_bucket" {
  value       = module.ml_platform.s3_bucket_names["models"]
  description = "S3 bucket for model storage"
}
\end{lstlisting}

\section{Kubernetes Deployment with Helm Charts}

Helm is the package manager for Kubernetes, enabling templated deployment of complex applications.

\subsection{Helm Chart for ML Training Jobs}

\begin{lstlisting}[style=yaml, caption=Helm Chart for ML Training]
# charts/ml-training/Chart.yaml
apiVersion: v2
name: ml-training
description: Helm chart for ML training jobs
type: application
version: 1.0.0
appVersion: "1.0"

maintainers:
  - name: ML Platform Team
    email: ml-team@company.com

keywords:
  - machine-learning
  - training
  - pytorch
  - tensorflow
\end{lstlisting}

\begin{lstlisting}[style=yaml, caption=Training Job Values]
# charts/ml-training/values.yaml
# Default values for ml-training chart

replicaCount: 1

image:
  repository: ${ECR_REPOSITORY}/training
  pullPolicy: IfNotPresent
  tag: "latest"

imagePullSecrets: []
nameOverride: ""
fullnameOverride: ""

# Training configuration
training:
  # Job type: single-node or distributed
  jobType: single-node

  # Training framework: pytorch, tensorflow, scikit-learn
  framework: pytorch

  # Training script and arguments
  command:
    - "python"
    - "train.py"
  args:
    - "--epochs=100"
    - "--batch-size=32"
    - "--learning-rate=0.001"

  # Data paths
  dataPath: /data
  modelOutputPath: /models

  # Environment variables
  env:
    - name: MLFLOW_TRACKING_URI
      value: "http://mlflow-server:5000"
    - name: AWS_DEFAULT_REGION
      value: "us-west-2"

  # Secrets (for credentials)
  envFrom:
    - secretRef:
        name: ml-credentials

# Resource requirements
resources:
  limits:
    nvidia.com/gpu: 1
    memory: 32Gi
    cpu: 8
  requests:
    nvidia.com/gpu: 1
    memory: 16Gi
    cpu: 4

# Node selection
nodeSelector:
  workload: training
  gpu: nvidia-v100

tolerations:
  - key: nvidia.com/gpu
    operator: Equal
    value: "true"
    effect: NoSchedule

affinity: {}

# Persistent volume claims
persistence:
  enabled: true
  data:
    storageClass: "gp3"
    accessMode: ReadOnlyMany
    size: 100Gi
    mountPath: /data
  models:
    storageClass: "gp3"
    accessMode: ReadWriteOnce
    size: 50Gi
    mountPath: /models

# Service account
serviceAccount:
  create: true
  annotations:
    eks.amazonaws.com/role-arn: arn:aws:iam::ACCOUNT_ID:role/ml-training-role
  name: ml-training

# Security context
securityContext:
  runAsNonRoot: true
  runAsUser: 1000
  fsGroup: 1000
  capabilities:
    drop:
      - ALL

podSecurityContext:
  seccompProfile:
    type: RuntimeDefault

# Monitoring
monitoring:
  enabled: true
  serviceMonitor:
    enabled: true
    interval: 30s
    scrapeTimeout: 10s

# Logging
logging:
  level: INFO
  format: json
\end{lstlisting}

\begin{lstlisting}[style=yaml, caption=Helm Template for Training Job]
# charts/ml-training/templates/job.yaml
{{- if eq .Values.training.jobType "single-node" }}
apiVersion: batch/v1
kind: Job
metadata:
  name: {{ include "ml-training.fullname" . }}
  labels:
    {{- include "ml-training.labels" . | nindent 4 }}
spec:
  backoffLimit: 3
  completions: 1
  parallelism: 1
  template:
    metadata:
      annotations:
        sidecar.istio.io/inject: "false"
      labels:
        {{- include "ml-training.selectorLabels" . | nindent 8 }}
        training-job: "true"
    spec:
      restartPolicy: Never
      serviceAccountName: {{ include "ml-training.serviceAccountName" . }}

      {{- with .Values.imagePullSecrets }}
      imagePullSecrets:
        {{- toYaml . | nindent 8 }}
      {{- end }}

      {{- with .Values.nodeSelector }}
      nodeSelector:
        {{- toYaml . | nindent 8 }}
      {{- end }}

      {{- with .Values.tolerations }}
      tolerations:
        {{- toYaml . | nindent 8 }}
      {{- end }}

      {{- with .Values.affinity }}
      affinity:
        {{- toYaml . | nindent 8 }}
      {{- end }}

      volumes:
        {{- if .Values.persistence.enabled }}
        - name: data
          persistentVolumeClaim:
            claimName: {{ include "ml-training.fullname" . }}-data
        - name: models
          persistentVolumeClaim:
            claimName: {{ include "ml-training.fullname" . }}-models
        {{- end }}
        - name: shared-memory
          emptyDir:
            medium: Memory
            sizeLimit: 8Gi

      initContainers:
        - name: data-validation
          image: {{ .Values.image.repository }}:{{ .Values.image.tag }}
          command:
            - python
            - -c
            - |
              import os
              import sys
              data_path = "{{ .Values.training.dataPath }}"
              if not os.path.exists(data_path):
                  print(f"Data path {data_path} does not exist")
                  sys.exit(1)
              if not os.listdir(data_path):
                  print(f"Data path {data_path} is empty")
                  sys.exit(1)
              print(f"Data validation passed: {len(os.listdir(data_path))} files found")
          volumeMounts:
            - name: data
              mountPath: {{ .Values.training.dataPath }}

      containers:
        - name: training
          image: {{ .Values.image.repository }}:{{ .Values.image.tag }}
          imagePullPolicy: {{ .Values.image.pullPolicy }}

          command: {{ toYaml .Values.training.command | nindent 12 }}
          {{- with .Values.training.args }}
          args: {{ toYaml . | nindent 12 }}
          {{- end }}

          env:
            {{- range .Values.training.env }}
            - name: {{ .name }}
              value: {{ .value | quote }}
            {{- end }}
            - name: PYTHONUNBUFFERED
              value: "1"
            - name: POD_NAME
              valueFrom:
                fieldRef:
                  fieldPath: metadata.name
            - name: POD_NAMESPACE
              valueFrom:
                fieldRef:
                  fieldPath: metadata.namespace

          {{- with .Values.training.envFrom }}
          envFrom: {{ toYaml . | nindent 12 }}
          {{- end }}

          volumeMounts:
            - name: data
              mountPath: {{ .Values.training.dataPath }}
              readOnly: true
            - name: models
              mountPath: {{ .Values.training.modelOutputPath }}
            - name: shared-memory
              mountPath: /dev/shm

          resources: {{ toYaml .Values.resources | nindent 12 }}

          securityContext: {{ toYaml .Values.securityContext | nindent 12 }}

      securityContext: {{ toYaml .Values.podSecurityContext | nindent 8 }}
{{- end }}
\end{lstlisting}

\begin{lstlisting}[style=yaml, caption=Helm Chart for Model Serving]
# charts/ml-serving/values.yaml
replicaCount: 3

image:
  repository: ${ECR_REPOSITORY}/inference
  pullPolicy: IfNotPresent
  tag: "latest"

service:
  type: ClusterIP
  port: 8080
  targetPort: 8080
  annotations:
    service.beta.kubernetes.io/aws-load-balancer-type: "nlb"

ingress:
  enabled: true
  className: "nginx"
  annotations:
    cert-manager.io/cluster-issuer: "letsencrypt-prod"
    nginx.ingress.kubernetes.io/ssl-redirect: "true"
    nginx.ingress.kubernetes.io/force-ssl-redirect: "true"
  hosts:
    - host: api.ml-platform.com
      paths:
        - path: /predict
          pathType: Prefix
  tls:
    - secretName: ml-api-tls
      hosts:
        - api.ml-platform.com

# Model configuration
model:
  name: my-model
  version: production
  registry: mlflow
  registryUri: http://mlflow-server:5000

# Autoscaling
autoscaling:
  enabled: true
  minReplicas: 3
  maxReplicas: 50
  targetCPUUtilizationPercentage: 70
  targetMemoryUtilizationPercentage: 80

  # Custom metrics for ML workloads
  metrics:
    - type: Pods
      pods:
        metric:
          name: inference_queue_depth
        target:
          type: AverageValue
          averageValue: "10"
    - type: Pods
      pods:
        metric:
          name: inference_latency_p95
        target:
          type: AverageValue
          averageValue: "100m"

resources:
  limits:
    memory: 8Gi
    cpu: 4
  requests:
    memory: 4Gi
    cpu: 2

# Health checks
livenessProbe:
  httpGet:
    path: /health
    port: 8080
  initialDelaySeconds: 30
  periodSeconds: 10
  timeoutSeconds: 5
  failureThreshold: 3

readinessProbe:
  httpGet:
    path: /ready
    port: 8080
  initialDelaySeconds: 15
  periodSeconds: 5
  timeoutSeconds: 3
  failureThreshold: 3

# Pod disruption budget
podDisruptionBudget:
  enabled: true
  minAvailable: 2
\end{lstlisting}

\subsection{Distributed Training with Kubernetes}

\begin{lstlisting}[style=yaml, caption=PyTorch Distributed Training Job]
# charts/ml-training/templates/pytorch-distributed.yaml
{{- if eq .Values.training.jobType "distributed" }}
apiVersion: kubeflow.org/v1
kind: PyTorchJob
metadata:
  name: {{ include "ml-training.fullname" . }}
  labels:
    {{- include "ml-training.labels" . | nindent 4 }}
spec:
  pytorchReplicaSpecs:
    Master:
      replicas: 1
      restartPolicy: OnFailure
      template:
        metadata:
          labels:
            {{- include "ml-training.selectorLabels" . | nindent 12 }}
            role: master
        spec:
          containers:
            - name: pytorch
              image: {{ .Values.image.repository }}:{{ .Values.image.tag }}
              command: {{ toYaml .Values.training.command | nindent 16 }}
              args:
                {{- range .Values.training.args }}
                - {{ . }}
                {{- end }}
                - "--distributed"
                - "--backend=nccl"
              resources:
                limits:
                  nvidia.com/gpu: {{ .Values.resources.limits."nvidia.com/gpu" }}
                  memory: {{ .Values.resources.limits.memory }}
                  cpu: {{ .Values.resources.limits.cpu }}
                requests:
                  nvidia.com/gpu: {{ .Values.resources.requests."nvidia.com/gpu" }}
                  memory: {{ .Values.resources.requests.memory }}
                  cpu: {{ .Values.resources.requests.cpu }}
              volumeMounts:
                - name: data
                  mountPath: {{ .Values.training.dataPath }}
                - name: models
                  mountPath: {{ .Values.training.modelOutputPath }}
                - name: shared-memory
                  mountPath: /dev/shm
          volumes:
            - name: data
              persistentVolumeClaim:
                claimName: training-data
            - name: models
              persistentVolumeClaim:
                claimName: training-models
            - name: shared-memory
              emptyDir:
                medium: Memory
                sizeLimit: 8Gi
          nodeSelector: {{ toYaml .Values.nodeSelector | nindent 12 }}
          tolerations: {{ toYaml .Values.tolerations | nindent 12 }}

    Worker:
      replicas: {{ .Values.training.workers | default 3 }}
      restartPolicy: OnFailure
      template:
        metadata:
          labels:
            {{- include "ml-training.selectorLabels" . | nindent 12 }}
            role: worker
        spec:
          containers:
            - name: pytorch
              image: {{ .Values.image.repository }}:{{ .Values.image.tag }}
              command: {{ toYaml .Values.training.command | nindent 16 }}
              args:
                {{- range .Values.training.args }}
                - {{ . }}
                {{- end }}
                - "--distributed"
                - "--backend=nccl"
              resources:
                limits:
                  nvidia.com/gpu: {{ .Values.resources.limits."nvidia.com/gpu" }}
                  memory: {{ .Values.resources.limits.memory }}
                  cpu: {{ .Values.resources.limits.cpu }}
                requests:
                  nvidia.com/gpu: {{ .Values.resources.requests."nvidia.com/gpu" }}
                  memory: {{ .Values.resources.requests.memory }}
                  cpu: {{ .Values.resources.requests.cpu }}
              volumeMounts:
                - name: data
                  mountPath: {{ .Values.training.dataPath }}
                  readOnly: true
                - name: models
                  mountPath: {{ .Values.training.modelOutputPath }}
                - name: shared-memory
                  mountPath: /dev/shm
          volumes:
            - name: data
              persistentVolumeClaim:
                claimName: training-data
            - name: models
              persistentVolumeClaim:
                claimName: training-models
            - name: shared-memory
              emptyDir:
                medium: Memory
                sizeLimit: 8Gi
          nodeSelector: {{ toYaml .Values.nodeSelector | nindent 12 }}
          tolerations: {{ toYaml .Values.tolerations | nindent 12 }}
{{- end }}
\end{lstlisting}

\section{Multi-Cloud Resource Management}

Modern ML platforms often span multiple cloud providers for redundancy, cost optimization, and feature availability.

\subsection{Google Cloud Platform (GCP) Infrastructure}

\begin{lstlisting}[language=Python, caption=Terraform Configuration for GCP ML Platform]
# modules/ml-platform-gcp/main.tf
"""
GCP ML Platform infrastructure module
Provides GKE, Cloud Storage, Vertex AI integration
"""

terraform {
  required_providers {
    google = {
      source  = "hashicorp/google"
      version = "~> 5.0"
    }
    google-beta = {
      source  = "hashicorp/google-beta"
      version = "~> 5.0"
    }
  }
}

provider "google" {
  project = var.gcp_project_id
  region  = var.gcp_region
}

provider "google-beta" {
  project = var.gcp_project_id
  region  = var.gcp_region
}

locals {
  name_prefix = "${var.project_name}-${var.environment}"
}

# VPC Network
resource "google_compute_network" "ml_vpc" {
  name                    = "${local.name_prefix}-vpc"
  auto_create_subnetworks = false
  routing_mode            = "REGIONAL"
}

resource "google_compute_subnetwork" "ml_subnet" {
  name          = "${local.name_prefix}-subnet"
  ip_cidr_range = "10.0.0.0/20"
  region        = var.gcp_region
  network       = google_compute_network.ml_vpc.id

  secondary_ip_range {
    range_name    = "pods"
    ip_cidr_range = "10.4.0.0/14"
  }

  secondary_ip_range {
    range_name    = "services"
    ip_cidr_range = "10.8.0.0/20"
  }

  private_ip_google_access = true
}

# Cloud NAT for private GKE nodes
resource "google_compute_router" "ml_router" {
  name    = "${local.name_prefix}-router"
  network = google_compute_network.ml_vpc.id
  region  = var.gcp_region
}

resource "google_compute_router_nat" "ml_nat" {
  name                               = "${local.name_prefix}-nat"
  router                             = google_compute_router.ml_router.name
  region                             = var.gcp_region
  nat_ip_allocate_option             = "AUTO_ONLY"
  source_subnetwork_ip_ranges_to_nat = "ALL_SUBNETWORKS_ALL_IP_RANGES"
}

# GKE Cluster
resource "google_container_cluster" "ml_cluster" {
  name     = "${local.name_prefix}-gke"
  location = var.gcp_region

  network    = google_compute_network.ml_vpc.name
  subnetwork = google_compute_subnetwork.ml_subnet.name

  # Private cluster configuration
  private_cluster_config {
    enable_private_nodes    = true
    enable_private_endpoint = false
    master_ipv4_cidr_block  = "172.16.0.0/28"
  }

  ip_allocation_policy {
    cluster_secondary_range_name  = "pods"
    services_secondary_range_name = "services"
  }

  # Workload Identity
  workload_identity_config {
    workload_pool = "${var.gcp_project_id}.svc.id.goog"
  }

  # Addons
  addons_config {
    http_load_balancing {
      disabled = false
    }
    horizontal_pod_autoscaling {
      disabled = false
    }
    gce_persistent_disk_csi_driver_config {
      enabled = true
    }
  }

  # Maintenance window
  maintenance_policy {
    daily_maintenance_window {
      start_time = "03:00"
    }
  }

  # Remove default node pool
  remove_default_node_pool = true
  initial_node_count       = 1

  deletion_protection = var.environment == "production"
}

# CPU Node Pool
resource "google_container_node_pool" "cpu_nodes" {
  name       = "${local.name_prefix}-cpu-pool"
  location   = var.gcp_region
  cluster    = google_container_cluster.ml_cluster.name
  node_count = var.environment == "production" ? 3 : 2

  autoscaling {
    min_node_count = var.environment == "production" ? 2 : 1
    max_node_count = 10
  }

  node_config {
    machine_type = "n2-standard-8"
    disk_size_gb = 100
    disk_type    = "pd-standard"

    oauth_scopes = [
      "https://www.googleapis.com/auth/cloud-platform"
    ]

    labels = {
      workload    = "general"
      environment = var.environment
    }

    tags = ["ml-workload", var.environment]

    workload_metadata_config {
      mode = "GKE_METADATA"
    }

    shielded_instance_config {
      enable_secure_boot          = true
      enable_integrity_monitoring = true
    }
  }

  management {
    auto_repair  = true
    auto_upgrade = true
  }
}

# GPU Node Pool for Training
resource "google_container_node_pool" "gpu_nodes" {
  name       = "${local.name_prefix}-gpu-pool"
  location   = var.gcp_region
  cluster    = google_container_cluster.ml_cluster.name
  node_count = 0  # Scale from zero

  autoscaling {
    min_node_count = 0
    max_node_count = 20
  }

  node_config {
    machine_type = "n1-standard-8"
    disk_size_gb = 100

    guest_accelerator {
      type  = "nvidia-tesla-t4"
      count = 1

      gpu_driver_installation_config {
        gpu_driver_version = "DEFAULT"
      }
    }

    oauth_scopes = [
      "https://www.googleapis.com/auth/cloud-platform"
    ]

    labels = {
      workload    = "training"
      gpu         = "nvidia-t4"
      environment = var.environment
    }

    taint {
      key    = "nvidia.com/gpu"
      value  = "true"
      effect = "NO_SCHEDULE"
    }

    tags = ["ml-gpu", var.environment]

    workload_metadata_config {
      mode = "GKE_METADATA"
    }
  }

  management {
    auto_repair  = true
    auto_upgrade = false  # Manual GPU node upgrades
  }
}

# Cloud Storage Buckets
resource "google_storage_bucket" "ml_data" {
  name          = "${local.name_prefix}-data"
  location      = var.gcp_region
  force_destroy = var.environment != "production"

  uniform_bucket_level_access = true

  versioning {
    enabled = true
  }

  lifecycle_rule {
    condition {
      age = 90
    }
    action {
      type          = "SetStorageClass"
      storage_class = "NEARLINE"
    }
  }

  lifecycle_rule {
    condition {
      age = 365
    }
    action {
      type          = "SetStorageClass"
      storage_class = "COLDLINE"
    }
  }
}

resource "google_storage_bucket" "ml_models" {
  name          = "${local.name_prefix}-models"
  location      = var.gcp_region
  force_destroy = var.environment != "production"

  uniform_bucket_level_access = true

  versioning {
    enabled = true
  }

  lifecycle_rule {
    condition {
      num_newer_versions = 5
    }
    action {
      type = "Delete"
    }
  }
}

# Artifact Registry for Container Images
resource "google_artifact_registry_repository" "ml_images" {
  location      = var.gcp_region
  repository_id = "${local.name_prefix}-images"
  format        = "DOCKER"

  cleanup_policies {
    id     = "delete-old-images"
    action = "DELETE"

    condition {
      tag_state  = "ANY"
      older_than = "2592000s"  # 30 days
    }
  }
}

# Cloud SQL for Metadata
resource "google_sql_database_instance" "ml_metadata" {
  name             = "${local.name_prefix}-metadata"
  database_version = "POSTGRES_15"
  region           = var.gcp_region

  deletion_protection = var.environment == "production"

  settings {
    tier              = var.environment == "production" ? "db-custom-4-16384" : "db-custom-2-7680"
    availability_type = var.environment == "production" ? "REGIONAL" : "ZONAL"
    disk_size         = 100
    disk_type         = "PD_SSD"
    disk_autoresize   = true

    backup_configuration {
      enabled                        = true
      start_time                     = "03:00"
      point_in_time_recovery_enabled = var.environment == "production"
      backup_retention_settings {
        retained_backups = var.environment == "production" ? 30 : 7
      }
    }

    ip_configuration {
      ipv4_enabled    = false
      private_network = google_compute_network.ml_vpc.id
    }

    insights_config {
      query_insights_enabled  = true
      query_string_length     = 1024
      record_application_tags = true
      record_client_address   = true
    }
  }

  depends_on = [google_service_networking_connection.private_vpc_connection]
}

# Private Service Connection
resource "google_compute_global_address" "private_ip_address" {
  name          = "${local.name_prefix}-private-ip"
  purpose       = "VPC_PEERING"
  address_type  = "INTERNAL"
  prefix_length = 16
  network       = google_compute_network.ml_vpc.id
}

resource "google_service_networking_connection" "private_vpc_connection" {
  network                 = google_compute_network.ml_vpc.id
  service                 = "servicenetworking.googleapis.com"
  reserved_peering_ranges = [google_compute_global_address.private_ip_address.name]
}

# Memorystore Redis
resource "google_redis_instance" "ml_cache" {
  name           = "${local.name_prefix}-redis"
  tier           = var.environment == "production" ? "STANDARD_HA" : "BASIC"
  memory_size_gb = var.environment == "production" ? 5 : 1
  region         = var.gcp_region

  redis_version     = "REDIS_7_0"
  display_name      = "${local.name_prefix} ML Cache"
  authorized_network = google_compute_network.ml_vpc.id

  transit_encryption_mode = "SERVER_AUTHENTICATION"
  auth_enabled            = true

  maintenance_policy {
    weekly_maintenance_window {
      day = "SUNDAY"
      start_time {
        hours   = 3
        minutes = 0
      }
    }
  }
}

# Vertex AI Workbench Instance
resource "google_workbench_instance" "ml_notebook" {
  name     = "${local.name_prefix}-notebook"
  location = "${var.gcp_region}-a"

  gce_setup {
    machine_type = "n1-standard-4"

    accelerator_configs {
      type       = "NVIDIA_TESLA_T4"
      core_count = 1
    }

    vm_image {
      project = "deeplearning-platform-release"
      family  = "common-cu113-notebooks"
    }

    network_interfaces {
      network = google_compute_network.ml_vpc.id
      subnet  = google_compute_subnetwork.ml_subnet.id
    }

    boot_disk {
      disk_size_gb = 100
      disk_type    = "PD_SSD"
    }

    data_disks {
      disk_size_gb = 200
      disk_type    = "PD_STANDARD"
    }

    metadata = {
      terraform = "true"
    }
  }
}

# Outputs
output "gke_cluster_name" {
  value = google_container_cluster.ml_cluster.name
}

output "gke_cluster_endpoint" {
  value     = google_container_cluster.ml_cluster.endpoint
  sensitive = true
}

output "gcs_data_bucket" {
  value = google_storage_bucket.ml_data.name
}

output "gcs_models_bucket" {
  value = google_storage_bucket.ml_models.name
}

output "artifact_registry_url" {
  value = "${var.gcp_region}-docker.pkg.dev/${var.gcp_project_id}/${google_artifact_registry_repository.ml_images.repository_id}"
}
\end{lstlisting}

\subsection{Azure Machine Learning Infrastructure}

\begin{lstlisting}[language=Python, caption=Terraform Configuration for Azure ML Platform]
# modules/ml-platform-azure/main.tf
"""
Azure ML Platform infrastructure module
Provides AKS, Azure ML, Blob Storage
"""

terraform {
  required_providers {
    azurerm = {
      source  = "hashicorp/azurerm"
      version = "~> 3.0"
    }
  }
}

provider "azurerm" {
  features {
    key_vault {
      purge_soft_delete_on_destroy = var.environment != "production"
    }
    resource_group {
      prevent_deletion_if_contains_resources = var.environment == "production"
    }
  }
}

locals {
  name_prefix = "${var.project_name}-${var.environment}"
  location    = var.azure_location
}

# Resource Group
resource "azurerm_resource_group" "ml_rg" {
  name     = "${local.name_prefix}-rg"
  location = local.location

  tags = {
    Environment = var.environment
    Project     = var.project_name
    ManagedBy   = "Terraform"
  }
}

# Virtual Network
resource "azurerm_virtual_network" "ml_vnet" {
  name                = "${local.name_prefix}-vnet"
  location            = azurerm_resource_group.ml_rg.location
  resource_group_name = azurerm_resource_group.ml_rg.name
  address_space       = ["10.0.0.0/16"]

  tags = azurerm_resource_group.ml_rg.tags
}

resource "azurerm_subnet" "aks_subnet" {
  name                 = "${local.name_prefix}-aks-subnet"
  resource_group_name  = azurerm_resource_group.ml_rg.name
  virtual_network_name = azurerm_virtual_network.ml_vnet.name
  address_prefixes     = ["10.0.1.0/24"]
}

resource "azurerm_subnet" "ml_subnet" {
  name                 = "${local.name_prefix}-ml-subnet"
  resource_group_name  = azurerm_resource_group.ml_rg.name
  virtual_network_name = azurerm_virtual_network.ml_vnet.name
  address_prefixes     = ["10.0.2.0/24"]

  delegation {
    name = "ml-delegation"

    service_delegation {
      name = "Microsoft.MachineLearningServices/workspaces"
      actions = [
        "Microsoft.Network/virtualNetworks/subnets/join/action",
      ]
    }
  }
}

# AKS Cluster
resource "azurerm_kubernetes_cluster" "ml_aks" {
  name                = "${local.name_prefix}-aks"
  location            = azurerm_resource_group.ml_rg.location
  resource_group_name = azurerm_resource_group.ml_rg.name
  dns_prefix          = "${local.name_prefix}-aks"

  kubernetes_version        = var.kubernetes_version
  automatic_channel_upgrade = "patch"
  sku_tier                  = var.environment == "production" ? "Standard" : "Free"

  default_node_pool {
    name                = "default"
    node_count          = var.environment == "production" ? 3 : 2
    vm_size             = "Standard_D4s_v3"
    vnet_subnet_id      = azurerm_subnet.aks_subnet.id
    enable_auto_scaling = true
    min_count           = var.environment == "production" ? 2 : 1
    max_count           = 10

    upgrade_settings {
      max_surge = "10%"
    }

    tags = azurerm_resource_group.ml_rg.tags
  }

  identity {
    type = "SystemAssigned"
  }

  network_profile {
    network_plugin    = "azure"
    network_policy    = "azure"
    load_balancer_sku = "standard"
    service_cidr      = "10.1.0.0/16"
    dns_service_ip    = "10.1.0.10"
  }

  oms_agent {
    log_analytics_workspace_id = azurerm_log_analytics_workspace.ml_logs.id
  }

  azure_active_directory_role_based_access_control {
    managed                = true
    azure_rbac_enabled     = true
    admin_group_object_ids = var.admin_group_object_ids
  }

  tags = azurerm_resource_group.ml_rg.tags
}

# GPU Node Pool
resource "azurerm_kubernetes_cluster_node_pool" "gpu_pool" {
  name                  = "gpupool"
  kubernetes_cluster_id = azurerm_kubernetes_cluster.ml_aks.id
  vm_size               = "Standard_NC6s_v3"
  node_count            = 0
  enable_auto_scaling   = true
  min_count             = 0
  max_count             = 10
  vnet_subnet_id        = azurerm_subnet.aks_subnet.id

  node_labels = {
    workload = "training"
    gpu      = "nvidia-v100"
  }

  node_taints = [
    "nvidia.com/gpu=true:NoSchedule"
  ]

  tags = azurerm_resource_group.ml_rg.tags
}

# Storage Account
resource "azurerm_storage_account" "ml_storage" {
  name                     = replace("${local.name_prefix}storage", "-", "")
  resource_group_name      = azurerm_resource_group.ml_rg.name
  location                 = azurerm_resource_group.ml_rg.location
  account_tier             = "Standard"
  account_replication_type = var.environment == "production" ? "GRS" : "LRS"
  account_kind             = "StorageV2"

  is_hns_enabled = true  # Enable hierarchical namespace for Data Lake

  network_rules {
    default_action             = "Deny"
    ip_rules                   = []
    virtual_network_subnet_ids = [azurerm_subnet.ml_subnet.id]
    bypass                     = ["AzureServices"]
  }

  blob_properties {
    versioning_enabled = true

    delete_retention_policy {
      days = var.environment == "production" ? 30 : 7
    }

    container_delete_retention_policy {
      days = var.environment == "production" ? 30 : 7
    }
  }

  tags = azurerm_resource_group.ml_rg.tags
}

resource "azurerm_storage_container" "ml_data" {
  name                  = "ml-data"
  storage_account_name  = azurerm_storage_account.ml_storage.name
  container_access_type = "private"
}

resource "azurerm_storage_container" "ml_models" {
  name                  = "ml-models"
  storage_account_name  = azurerm_storage_account.ml_storage.name
  container_access_type = "private"
}

# Container Registry
resource "azurerm_container_registry" "ml_acr" {
  name                = replace("${local.name_prefix}acr", "-", "")
  resource_group_name = azurerm_resource_group.ml_rg.name
  location            = azurerm_resource_group.ml_rg.location
  sku                 = var.environment == "production" ? "Premium" : "Standard"
  admin_enabled       = false

  identity {
    type = "SystemAssigned"
  }

  georeplications = var.environment == "production" ? [
    {
      location = var.geo_replication_location
      tags     = {}
    }
  ] : []

  network_rule_set {
    default_action = "Deny"

    ip_rule {
      action   = "Allow"
      ip_range = azurerm_virtual_network.ml_vnet.address_space[0]
    }
  }

  retention_policy {
    days    = 7
    enabled = true
  }

  trust_policy {
    enabled = true
  }

  tags = azurerm_resource_group.ml_rg.tags
}

# Key Vault
resource "azurerm_key_vault" "ml_kv" {
  name                = "${local.name_prefix}-kv"
  location            = azurerm_resource_group.ml_rg.location
  resource_group_name = azurerm_resource_group.ml_rg.name
  tenant_id           = data.azurerm_client_config.current.tenant_id
  sku_name            = var.environment == "production" ? "premium" : "standard"

  purge_protection_enabled   = var.environment == "production"
  soft_delete_retention_days = var.environment == "production" ? 90 : 7

  network_acls {
    default_action             = "Deny"
    bypass                     = "AzureServices"
    virtual_network_subnet_ids = [azurerm_subnet.ml_subnet.id]
  }

  tags = azurerm_resource_group.ml_rg.tags
}

# Azure ML Workspace
resource "azurerm_machine_learning_workspace" "ml_workspace" {
  name                    = "${local.name_prefix}-mlworkspace"
  location                = azurerm_resource_group.ml_rg.location
  resource_group_name     = azurerm_resource_group.ml_rg.name
  application_insights_id = azurerm_application_insights.ml_insights.id
  key_vault_id            = azurerm_key_vault.ml_kv.id
  storage_account_id      = azurerm_storage_account.ml_storage.id
  container_registry_id   = azurerm_container_registry.ml_acr.id

  identity {
    type = "SystemAssigned"
  }

  public_network_access_enabled = var.environment != "production"

  tags = azurerm_resource_group.ml_rg.tags
}

# Log Analytics Workspace
resource "azurerm_log_analytics_workspace" "ml_logs" {
  name                = "${local.name_prefix}-logs"
  location            = azurerm_resource_group.ml_rg.location
  resource_group_name = azurerm_resource_group.ml_rg.name
  sku                 = "PerGB2018"
  retention_in_days   = var.environment == "production" ? 30 : 7

  tags = azurerm_resource_group.ml_rg.tags
}

# Application Insights
resource "azurerm_application_insights" "ml_insights" {
  name                = "${local.name_prefix}-insights"
  location            = azurerm_resource_group.ml_rg.location
  resource_group_name = azurerm_resource_group.ml_rg.name
  workspace_id        = azurerm_log_analytics_workspace.ml_logs.id
  application_type    = "other"

  tags = azurerm_resource_group.ml_rg.tags
}

# PostgreSQL Flexible Server
resource "azurerm_postgresql_flexible_server" "ml_db" {
  name                   = "${local.name_prefix}-psql"
  resource_group_name    = azurerm_resource_group.ml_rg.name
  location               = azurerm_resource_group.ml_rg.location
  version                = "15"
  administrator_login    = var.db_admin_username
  administrator_password = var.db_admin_password
  zone                   = "1"

  storage_mb = 32768

  sku_name = var.environment == "production" ? "GP_Standard_D4s_v3" : "GP_Standard_D2s_v3"

  backup_retention_days        = var.environment == "production" ? 30 : 7
  geo_redundant_backup_enabled = var.environment == "production"

  high_availability {
    mode                      = var.environment == "production" ? "ZoneRedundant" : "Disabled"
    standby_availability_zone = var.environment == "production" ? "2" : null
  }

  tags = azurerm_resource_group.ml_rg.tags
}

data "azurerm_client_config" "current" {}

# Outputs
output "aks_cluster_name" {
  value = azurerm_kubernetes_cluster.ml_aks.name
}

output "aks_cluster_id" {
  value = azurerm_kubernetes_cluster.ml_aks.id
}

output "ml_workspace_id" {
  value = azurerm_machine_learning_workspace.ml_workspace.id
}

output "storage_account_name" {
  value = azurerm_storage_account.ml_storage.name
}

output "container_registry_login_server" {
  value = azurerm_container_registry.ml_acr.login_server
}
\end{lstlisting}

\section{Container Orchestration for ML Workloads}

Container orchestration enables efficient management of ML workloads at scale, with dynamic resource allocation and fault tolerance.

\subsection{Kubernetes Operators for ML}

\begin{lstlisting}[language=Python, caption=Custom Resource Definition for ML Training]
# k8s/crds/mltraining-crd.yaml
"""
Custom Resource Definition for ML Training Jobs
Abstracts complexity and provides high-level ML workflow management
"""

apiVersion: apiextensions.k8s.io/v1
kind: CustomResourceDefinition
metadata:
  name: mltrainings.ml.platform.io
spec:
  group: ml.platform.io
  versions:
    - name: v1alpha1
      served: true
      storage: true
      schema:
        openAPIV3Schema:
          type: object
          properties:
            spec:
              type: object
              properties:
                framework:
                  type: string
                  enum: [pytorch, tensorflow, scikit-learn, xgboost]
                  description: ML framework to use
                modelName:
                  type: string
                  description: Name of the model to train
                dataSource:
                  type: object
                  properties:
                    type:
                      type: string
                      enum: [s3, gcs, azure-blob, pvc]
                    path:
                      type: string
                    credentials:
                      type: string
                      description: Secret name for credentials
                trainingScript:
                  type: object
                  properties:
                    image:
                      type: string
                    command:
                      type: array
                      items:
                        type: string
                    args:
                      type: array
                      items:
                        type: string
                hyperparameters:
                  type: object
                  additionalProperties:
                    type: string
                resources:
                  type: object
                  properties:
                    gpuType:
                      type: string
                      enum: [v100, a100, t4]
                    gpuCount:
                      type: integer
                      minimum: 0
                      maximum: 8
                    memory:
                      type: string
                    cpu:
                      type: string
                distributed:
                  type: object
                  properties:
                    enabled:
                      type: boolean
                    workers:
                      type: integer
                      minimum: 1
                      maximum: 100
                    backend:
                      type: string
                      enum: [nccl, gloo, mpi]
                outputModel:
                  type: object
                  properties:
                    registry:
                      type: string
                    registryUri:
                      type: string
                    version:
                      type: string
            status:
              type: object
              properties:
                phase:
                  type: string
                  enum: [Pending, Running, Succeeded, Failed]
                startTime:
                  type: string
                  format: date-time
                completionTime:
                  type: string
                  format: date-time
                conditions:
                  type: array
                  items:
                    type: object
                    properties:
                      type:
                        type: string
                      status:
                        type: string
                      reason:
                        type: string
                      message:
                        type: string
                metrics:
                  type: object
                  additionalProperties:
                    type: number
      subresources:
        status: {}
      additionalPrinterColumns:
        - name: Framework
          type: string
          jsonPath: .spec.framework
        - name: Phase
          type: string
          jsonPath: .status.phase
        - name: Age
          type: date
          jsonPath: .metadata.creationTimestamp
  scope: Namespaced
  names:
    plural: mltrainings
    singular: mltraining
    kind: MLTraining
    shortNames:
      - mlt
\end{lstlisting}

\begin{lstlisting}[style=yaml, caption=Example ML Training Custom Resource]
# examples/training-job.yaml
apiVersion: ml.platform.io/v1alpha1
kind: MLTraining
metadata:
  name: resnet50-imagenet
  namespace: ml-workloads
spec:
  framework: pytorch
  modelName: resnet50-imagenet

  dataSource:
    type: s3
    path: s3://ml-data/imagenet
    credentials: aws-credentials

  trainingScript:
    image: 123456789.dkr.ecr.us-west-2.amazonaws.com/training:latest
    command: ["python"]
    args:
      - "train.py"
      - "--model=resnet50"
      - "--epochs=90"
      - "--batch-size=256"

  hyperparameters:
    learning_rate: "0.1"
    weight_decay: "0.0001"
    momentum: "0.9"
    lr_schedule: "cosine"

  resources:
    gpuType: v100
    gpuCount: 4
    memory: "128Gi"
    cpu: "32"

  distributed:
    enabled: true
    workers: 4
    backend: nccl

  outputModel:
    registry: mlflow
    registryUri: http://mlflow-server:5000
    version: "1.0.0"
\end{lstlisting}

\subsection{Autoscaling ML Workloads}

\begin{lstlisting}[style=yaml, caption=Kubernetes Cluster Autoscaler for ML]
# k8s/autoscaling/cluster-autoscaler.yaml
apiVersion: v1
kind: ConfigMap
metadata:
  name: cluster-autoscaler-priority-expander
  namespace: kube-system
data:
  priorities: |
    10:
      - .*-gpu-training.*  # Highest priority for GPU training nodes
    5:
      - .*-inference.*     # Medium priority for inference nodes
    1:
      - .*-general.*       # Lowest priority for general workloads
---
apiVersion: apps/v1
kind: Deployment
metadata:
  name: cluster-autoscaler
  namespace: kube-system
  labels:
    app: cluster-autoscaler
spec:
  replicas: 1
  selector:
    matchLabels:
      app: cluster-autoscaler
  template:
    metadata:
      labels:
        app: cluster-autoscaler
    spec:
      serviceAccountName: cluster-autoscaler
      containers:
        - name: cluster-autoscaler
          image: registry.k8s.io/autoscaling/cluster-autoscaler:v1.28.0
          command:
            - ./cluster-autoscaler
            - --v=4
            - --stderrthreshold=info
            - --cloud-provider=aws
            - --skip-nodes-with-local-storage=false
            - --skip-nodes-with-system-pods=false
            - --expander=priority
            - --balance-similar-node-groups
            - --scale-down-delay-after-add=5m
            - --scale-down-unneeded-time=5m
            - --scale-down-utilization-threshold=0.5
            - --max-node-provision-time=15m
            - --node-group-auto-discovery=asg:tag=k8s.io/cluster-autoscaler/enabled,k8s.io/cluster-autoscaler/${CLUSTER_NAME}
          env:
            - name: CLUSTER_NAME
              valueFrom:
                configMapKeyRef:
                  name: cluster-info
                  key: cluster-name
            - name: AWS_REGION
              value: us-west-2
          resources:
            limits:
              cpu: 100m
              memory: 600Mi
            requests:
              cpu: 100m
              memory: 600Mi
          volumeMounts:
            - name: ssl-certs
              mountPath: /etc/ssl/certs/ca-certificates.crt
              readOnly: true
      volumes:
        - name: ssl-certs
          hostPath:
            path: /etc/ssl/certs/ca-bundle.crt
---
apiVersion: rbac.authorization.k8s.io/v1
kind: ClusterRole
metadata:
  name: cluster-autoscaler
rules:
  - apiGroups: [""]
    resources: ["events", "endpoints"]
    verbs: ["create", "patch"]
  - apiGroups: [""]
    resources: ["pods/eviction"]
    verbs: ["create"]
  - apiGroups: [""]
    resources: ["pods/status"]
    verbs: ["update"]
  - apiGroups: [""]
    resources: ["endpoints"]
    resourceNames: ["cluster-autoscaler"]
    verbs: ["get", "update"]
  - apiGroups: [""]
    resources: ["nodes"]
    verbs: ["watch", "list", "get", "update"]
  - apiGroups: [""]
    resources: ["pods", "services", "replicationcontrollers", "persistentvolumeclaims", "persistentvolumes"]
    verbs: ["watch", "list", "get"]
  - apiGroups: ["apps"]
    resources: ["statefulsets", "replicasets", "daemonsets"]
    verbs: ["watch", "list", "get"]
  - apiGroups: ["policy"]
    resources: ["poddisruptionbudgets"]
    verbs: ["watch", "list"]
  - apiGroups: ["storage.k8s.io"]
    resources: ["storageclasses", "csinodes", "csidrivers", "csistoragecapacities"]
    verbs: ["watch", "list", "get"]
  - apiGroups: ["batch"]
    resources: ["jobs"]
    verbs: ["watch", "list", "get"]
  - apiGroups: ["coordination.k8s.io"]
    resources: ["leases"]
    verbs: ["create"]
  - apiGroups: ["coordination.k8s.io"]
    resourceNames: ["cluster-autoscaler"]
    resources: ["leases"]
    verbs: ["get", "update"]
\end{lstlisting}

\subsection{Resource Quotas and Limits}

\begin{lstlisting}[style=yaml, caption=Resource Quotas for ML Namespaces]
# k8s/resource-management/quotas.yaml
apiVersion: v1
kind: ResourceQuota
metadata:
  name: ml-training-quota
  namespace: ml-training
spec:
  hard:
    requests.cpu: "200"
    requests.memory: 1Ti
    requests.nvidia.com/gpu: "32"
    limits.cpu: "400"
    limits.memory: 2Ti
    limits.nvidia.com/gpu: "64"
    persistentvolumeclaims: "20"
    requests.storage: 10Ti
  scopeSelector:
    matchExpressions:
      - operator: In
        scopeName: PriorityClass
        values: ["high", "medium"]
---
apiVersion: v1
kind: ResourceQuota
metadata:
  name: ml-inference-quota
  namespace: ml-inference
spec:
  hard:
    requests.cpu: "100"
    requests.memory: 512Gi
    limits.cpu: "200"
    limits.memory: 1Ti
    services.loadbalancers: "5"
    pods: "100"
---
apiVersion: v1
kind: LimitRange
metadata:
  name: ml-training-limits
  namespace: ml-training
spec:
  limits:
    - type: Pod
      max:
        cpu: "64"
        memory: "512Gi"
        nvidia.com/gpu: "8"
      min:
        cpu: "1"
        memory: "4Gi"
    - type: Container
      max:
        cpu: "64"
        memory: "512Gi"
        nvidia.com/gpu: "8"
      min:
        cpu: "100m"
        memory: "128Mi"
      default:
        cpu: "4"
        memory: "16Gi"
      defaultRequest:
        cpu: "2"
        memory: "8Gi"
    - type: PersistentVolumeClaim
      max:
        storage: "10Ti"
      min:
        storage: "10Gi"
\end{lstlisting}

\section{Disaster Recovery and Backup Strategies}

Robust backup and recovery strategies ensure ML system resilience against failures and data loss.

\subsection{Backup Automation with Terraform}

\begin{lstlisting}[language=Python, caption=Automated Backup Configuration]
# modules/backup-recovery/main.tf
"""
Automated backup and disaster recovery infrastructure
"""

terraform {
  required_providers {
    aws = {
      source  = "hashicorp/aws"
      version = "~> 5.0"
    }
  }
}

locals {
  backup_tags = {
    Backup      = "true"
    Environment = var.environment
    RetentionDays = var.backup_retention_days
  }
}

# AWS Backup Vault
resource "aws_backup_vault" "ml_backup_vault" {
  name        = "${var.project_name}-${var.environment}-vault"
  kms_key_arn = var.kms_key_arn

  tags = local.backup_tags
}

# Backup Plan
resource "aws_backup_plan" "ml_backup_plan" {
  name = "${var.project_name}-${var.environment}-backup-plan"

  # Daily backups
  rule {
    rule_name         = "daily_backup"
    target_vault_name = aws_backup_vault.ml_backup_vault.name
    schedule          = "cron(0 3 * * ? *)"  # 3 AM daily

    lifecycle {
      cold_storage_after = var.environment == "production" ? 30 : null
      delete_after       = var.environment == "production" ? 90 : 14
    }

    recovery_point_tags = local.backup_tags
  }

  # Weekly backups (longer retention)
  rule {
    rule_name         = "weekly_backup"
    target_vault_name = aws_backup_vault.ml_backup_vault.name
    schedule          = "cron(0 3 ? * 1 *)"  # Monday 3 AM

    lifecycle {
      cold_storage_after = var.environment == "production" ? 60 : null
      delete_after       = var.environment == "production" ? 365 : 30
    }

    recovery_point_tags = merge(local.backup_tags, {
      BackupType = "weekly"
    })
  }

  # Monthly backups (longest retention)
  rule {
    rule_name         = "monthly_backup"
    target_vault_name = aws_backup_vault.ml_backup_vault.name
    schedule          = "cron(0 3 1 * ? *)"  # First day of month

    lifecycle {
      cold_storage_after = var.environment == "production" ? 90 : null
      delete_after       = var.environment == "production" ? 2555 : 90  # ~7 years for production
    }

    recovery_point_tags = merge(local.backup_tags, {
      BackupType = "monthly"
    })
  }

  tags = local.backup_tags
}

# Backup Selection for RDS
resource "aws_backup_selection" "rds_backup" {
  name         = "rds-backup-selection"
  iam_role_arn = aws_iam_role.backup_role.arn
  plan_id      = aws_backup_plan.ml_backup_plan.id

  selection_tag {
    type  = "STRINGEQUALS"
    key   = "Backup"
    value = "true"
  }

  resources = var.rds_arns
}

# Backup Selection for EBS Volumes
resource "aws_backup_selection" "ebs_backup" {
  name         = "ebs-backup-selection"
  iam_role_arn = aws_iam_role.backup_role.arn
  plan_id      = aws_backup_plan.ml_backup_plan.id

  selection_tag {
    type  = "STRINGEQUALS"
    key   = "Backup"
    value = "true"
  }

  resources = var.ebs_volume_arns
}

# IAM Role for AWS Backup
resource "aws_iam_role" "backup_role" {
  name = "${var.project_name}-${var.environment}-backup-role"

  assume_role_policy = jsonencode({
    Version = "2012-10-17"
    Statement = [
      {
        Action = "sts:AssumeRole"
        Effect = "Allow"
        Principal = {
          Service = "backup.amazonaws.com"
        }
      }
    ]
  })

  tags = local.backup_tags
}

resource "aws_iam_role_policy_attachment" "backup_policy" {
  role       = aws_iam_role.backup_role.name
  policy_arn = "arn:aws:iam::aws:policy/service-role/AWSBackupServiceRolePolicyForBackup"
}

resource "aws_iam_role_policy_attachment" "restore_policy" {
  role       = aws_iam_role.backup_role.name
  policy_arn = "arn:aws:iam::aws:policy/service-role/AWSBackupServiceRolePolicyForRestores"
}

# S3 Bucket Replication for Models
resource "aws_s3_bucket_replication_configuration" "model_replication" {
  count = var.environment == "production" ? 1 : 0

  role   = aws_iam_role.replication_role[0].arn
  bucket = var.source_bucket_id

  rule {
    id     = "replicate-models"
    status = "Enabled"

    filter {
      prefix = "models/"
    }

    destination {
      bucket        = var.destination_bucket_arn
      storage_class = "STANDARD_IA"

      replication_time {
        status = "Enabled"
        time {
          minutes = 15
        }
      }

      metrics {
        status = "Enabled"
        event_threshold {
          minutes = 15
        }
      }
    }

    delete_marker_replication {
      status = "Enabled"
    }
  }
}

# IAM Role for S3 Replication
resource "aws_iam_role" "replication_role" {
  count = var.environment == "production" ? 1 : 0

  name = "${var.project_name}-${var.environment}-s3-replication-role"

  assume_role_policy = jsonencode({
    Version = "2012-10-17"
    Statement = [
      {
        Action = "sts:AssumeRole"
        Effect = "Allow"
        Principal = {
          Service = "s3.amazonaws.com"
        }
      }
    ]
  })

  tags = local.backup_tags
}

resource "aws_iam_role_policy" "replication_policy" {
  count = var.environment == "production" ? 1 : 0

  name = "${var.project_name}-s3-replication-policy"
  role = aws_iam_role.replication_role[0].id

  policy = jsonencode({
    Version = "2012-10-17"
    Statement = [
      {
        Effect = "Allow"
        Action = [
          "s3:GetReplicationConfiguration",
          "s3:ListBucket"
        ]
        Resource = var.source_bucket_arn
      },
      {
        Effect = "Allow"
        Action = [
          "s3:GetObjectVersionForReplication",
          "s3:GetObjectVersionAcl",
          "s3:GetObjectVersionTagging"
        ]
        Resource = "${var.source_bucket_arn}/*"
      },
      {
        Effect = "Allow"
        Action = [
          "s3:ReplicateObject",
          "s3:ReplicateDelete",
          "s3:ReplicateTags"
        ]
        Resource = "${var.destination_bucket_arn}/*"
      }
    ]
  })
}

# CloudWatch Alarms for Backup Failures
resource "aws_cloudwatch_metric_alarm" "backup_failure" {
  alarm_name          = "${var.project_name}-${var.environment}-backup-failure"
  comparison_operator = "GreaterThanThreshold"
  evaluation_periods  = 1
  metric_name         = "NumberOfBackupJobsFailed"
  namespace           = "AWS/Backup"
  period              = 3600
  statistic           = "Sum"
  threshold           = 0
  alarm_description   = "Alert when backup jobs fail"
  alarm_actions       = [var.sns_topic_arn]

  dimensions = {
    BackupVaultName = aws_backup_vault.ml_backup_vault.name
  }

  tags = local.backup_tags
}

# Outputs
output "backup_vault_arn" {
  value = aws_backup_vault.ml_backup_vault.arn
}

output "backup_plan_id" {
  value = aws_backup_plan.ml_backup_plan.id
}
\end{lstlisting}

\subsection{Disaster Recovery Procedures}

\begin{lstlisting}[language=Python, caption=Disaster Recovery Automation Script]
#!/usr/bin/env python3
"""
Disaster Recovery automation for ML infrastructure
Orchestrates recovery from backup vault
"""

import boto3
import logging
from datetime import datetime, timedelta
from typing import Dict, List, Optional
import argparse

logging.basicConfig(level=logging.INFO)
logger = logging.getLogger(__name__)

class DisasterRecoveryOrchestrator:
    """Orchestrates disaster recovery operations"""

    def __init__(self, region: str, environment: str):
        self.backup_client = boto3.client('backup', region_name=region)
        self.rds_client = boto3.client('rds', region_name=region)
        self.ec2_client = boto3.client('ec2', region_name=region)
        self.s3_client = boto3.client('s3', region_name=region)
        self.environment = environment
        self.region = region

    def list_recovery_points(
        self,
        backup_vault_name: str,
        resource_type: Optional[str] = None,
        created_after: Optional[datetime] = None
    ) -> List[Dict]:
        """List available recovery points"""

        params = {
            'BackupVaultName': backup_vault_name
        }

        if resource_type:
            params['ByResourceType'] = resource_type

        if created_after:
            params['ByCreatedAfter'] = created_after

        try:
            response = self.backup_client.list_recovery_points_by_backup_vault(
                **params
            )

            recovery_points = response.get('RecoveryPoints', [])
            logger.info(
                f"Found {len(recovery_points)} recovery points in {backup_vault_name}"
            )

            return recovery_points

        except Exception as e:
            logger.error(f"Failed to list recovery points: {e}")
            return []

    def select_latest_recovery_point(
        self,
        backup_vault_name: str,
        resource_arn: str
    ) -> Optional[str]:
        """Select the most recent recovery point for a resource"""

        try:
            response = self.backup_client.list_recovery_points_by_resource(
                ResourceArn=resource_arn
            )

            recovery_points = response.get('RecoveryPoints', [])

            # Filter completed recovery points
            completed = [
                rp for rp in recovery_points
                if rp['Status'] == 'COMPLETED'
            ]

            if not completed:
                logger.warning(f"No completed recovery points for {resource_arn}")
                return None

            # Sort by creation time and get latest
            latest = sorted(
                completed,
                key=lambda x: x['CreationDate'],
                reverse=True
            )[0]

            logger.info(
                f"Selected recovery point {latest['RecoveryPointArn']} "
                f"from {latest['CreationDate']}"
            )

            return latest['RecoveryPointArn']

        except Exception as e:
            logger.error(f"Failed to select recovery point: {e}")
            return None

    def restore_rds_instance(
        self,
        recovery_point_arn: str,
        target_identifier: str,
        subnet_group_name: str,
        security_group_ids: List[str]
    ) -> Optional[str]:
        """Restore RDS instance from backup"""

        logger.info(f"Starting RDS restore to {target_identifier}")

        try:
            # Start restore job
            response = self.backup_client.start_restore_job(
                RecoveryPointArn=recovery_point_arn,
                Metadata={
                    'DBInstanceIdentifier': target_identifier,
                    'DBSubnetGroupName': subnet_group_name,
                    'VpcSecurityGroupIds': ','.join(security_group_ids),
                    'PubliclyAccessible': 'false',
                    'MultiAZ': 'true' if self.environment == 'production' else 'false',
                    'StorageEncrypted': 'true'
                },
                IamRoleArn=self._get_restore_role_arn()
            )

            restore_job_id = response['RestoreJobId']
            logger.info(f"RDS restore job started: {restore_job_id}")

            return restore_job_id

        except Exception as e:
            logger.error(f"Failed to restore RDS instance: {e}")
            return None

    def restore_ebs_volume(
        self,
        recovery_point_arn: str,
        availability_zone: str,
        volume_type: str = 'gp3',
        iops: int = 3000
    ) -> Optional[str]:
        """Restore EBS volume from backup"""

        logger.info(f"Starting EBS volume restore in {availability_zone}")

        try:
            response = self.backup_client.start_restore_job(
                RecoveryPointArn=recovery_point_arn,
                Metadata={
                    'AvailabilityZone': availability_zone,
                    'VolumeType': volume_type,
                    'Iops': str(iops),
                    'Encrypted': 'true'
                },
                IamRoleArn=self._get_restore_role_arn()
            )

            restore_job_id = response['RestoreJobId']
            logger.info(f"EBS restore job started: {restore_job_id}")

            return restore_job_id

        except Exception as e:
            logger.error(f"Failed to restore EBS volume: {e}")
            return None

    def monitor_restore_job(self, restore_job_id: str) -> Dict:
        """Monitor restore job progress"""

        try:
            response = self.backup_client.describe_restore_job(
                RestoreJobId=restore_job_id
            )

            status = response['Status']
            logger.info(f"Restore job {restore_job_id} status: {status}")

            return response

        except Exception as e:
            logger.error(f"Failed to monitor restore job: {e}")
            return {}

    def replicate_s3_data(
        self,
        source_bucket: str,
        destination_bucket: str,
        prefix: str = ""
    ) -> bool:
        """Manually replicate S3 data for disaster recovery"""

        logger.info(
            f"Replicating from {source_bucket}/{prefix} "
            f"to {destination_bucket}/{prefix}"
        )

        try:
            paginator = self.s3_client.get_paginator('list_objects_v2')
            pages = paginator.paginate(Bucket=source_bucket, Prefix=prefix)

            objects_replicated = 0

            for page in pages:
                if 'Contents' not in page:
                    continue

                for obj in page['Contents']:
                    source_key = obj['Key']

                    # Copy object to destination
                    copy_source = {
                        'Bucket': source_bucket,
                        'Key': source_key
                    }

                    self.s3_client.copy_object(
                        CopySource=copy_source,
                        Bucket=destination_bucket,
                        Key=source_key
                    )

                    objects_replicated += 1

                    if objects_replicated % 100 == 0:
                        logger.info(f"Replicated {objects_replicated} objects")

            logger.info(
                f"Successfully replicated {objects_replicated} objects"
            )
            return True

        except Exception as e:
            logger.error(f"Failed to replicate S3 data: {e}")
            return False

    def _get_restore_role_arn(self) -> str:
        """Get IAM role ARN for restore operations"""
        # Construct ARN from account and environment
        account_id = boto3.client('sts').get_caller_identity()['Account']
        return (
            f"arn:aws:iam::{account_id}:role/"
            f"ml-platform-{self.environment}-backup-role"
        )

    def execute_full_recovery(
        self,
        backup_vault_name: str,
        recovery_config: Dict
    ) -> Dict:
        """Execute full disaster recovery procedure"""

        logger.info("Starting full disaster recovery procedure")

        results = {
            'rds_restores': [],
            'ebs_restores': [],
            's3_replications': [],
            'status': 'in_progress'
        }

        try:
            # Restore RDS instances
            for rds_config in recovery_config.get('rds_instances', []):
                recovery_point = self.select_latest_recovery_point(
                    backup_vault_name,
                    rds_config['source_arn']
                )

                if recovery_point:
                    restore_job = self.restore_rds_instance(
                        recovery_point,
                        rds_config['target_identifier'],
                        rds_config['subnet_group'],
                        rds_config['security_groups']
                    )

                    if restore_job:
                        results['rds_restores'].append({
                            'target': rds_config['target_identifier'],
                            'job_id': restore_job,
                            'status': 'started'
                        })

            # Restore EBS volumes
            for ebs_config in recovery_config.get('ebs_volumes', []):
                recovery_point = self.select_latest_recovery_point(
                    backup_vault_name,
                    ebs_config['source_arn']
                )

                if recovery_point:
                    restore_job = self.restore_ebs_volume(
                        recovery_point,
                        ebs_config['availability_zone'],
                        ebs_config.get('volume_type', 'gp3')
                    )

                    if restore_job:
                        results['ebs_restores'].append({
                            'az': ebs_config['availability_zone'],
                            'job_id': restore_job,
                            'status': 'started'
                        })

            # Replicate S3 data
            for s3_config in recovery_config.get('s3_replications', []):
                success = self.replicate_s3_data(
                    s3_config['source_bucket'],
                    s3_config['destination_bucket'],
                    s3_config.get('prefix', '')
                )

                results['s3_replications'].append({
                    'source': s3_config['source_bucket'],
                    'destination': s3_config['destination_bucket'],
                    'status': 'completed' if success else 'failed'
                })

            results['status'] = 'completed'
            logger.info("Disaster recovery procedure completed")

        except Exception as e:
            logger.error(f"Disaster recovery failed: {e}")
            results['status'] = 'failed'
            results['error'] = str(e)

        return results

def main():
    parser = argparse.ArgumentParser(
        description='ML Infrastructure Disaster Recovery'
    )
    parser.add_argument('--region', required=True, help='AWS region')
    parser.add_argument('--environment', required=True, help='Environment name')
    parser.add_argument('--vault', required=True, help='Backup vault name')
    parser.add_argument('--config', required=True, help='Recovery config file')

    args = parser.parse_args()

    # Load recovery configuration
    import json
    with open(args.config, 'r') as f:
        recovery_config = json.load(f)

    # Execute recovery
    orchestrator = DisasterRecoveryOrchestrator(args.region, args.environment)
    results = orchestrator.execute_full_recovery(args.vault, recovery_config)

    # Print results
    print(json.dumps(results, indent=2))

    return 0 if results['status'] == 'completed' else 1

if __name__ == '__main__':
    exit(main())
\end{lstlisting}

\section{Cost Optimization Through IaC}

Infrastructure as Code enables systematic cost optimization through automated resource management and intelligent provisioning strategies.

\subsection{Cost-Aware Resource Provisioning}

\begin{lstlisting}[language=Python, caption=Terraform Module for Cost-Optimized ML Infrastructure]
# modules/cost-optimization/main.tf
"""
Cost-optimized infrastructure module
Implements spot instances, auto-scaling, and resource lifecycle management
"""

terraform {
  required_providers {
    aws = {
      source  = "hashicorp/aws"
      version = "~> 5.0"
    }
  }
}

locals {
  cost_allocation_tags = {
    CostCenter  = var.cost_center
    Team        = var.team_name
    Environment = var.environment
    Project     = var.project_name
  }
}

# Spot Instance Fleet for Training
resource "aws_autoscaling_group" "training_spot_fleet" {
  name     = "${var.project_name}-training-spot"
  min_size = 0
  max_size = var.max_training_nodes
  desired_capacity = 0  # Scale from zero

  vpc_zone_identifier = var.subnet_ids

  mixed_instances_policy {
    instances_distribution {
      on_demand_base_capacity                  = 0
      on_demand_percentage_above_base_capacity = 20  # 80% spot
      spot_allocation_strategy                 = "price-capacity-optimized"
      spot_instance_pools                      = 4
    }

    launch_template {
      launch_template_specification {
        launch_template_id = aws_launch_template.training_spot.id
        version            = "$Latest"
      }

      # Multiple instance types for better spot availability
      override {
        instance_type     = "p3.2xlarge"
        weighted_capacity = 1
      }

      override {
        instance_type     = "p3.8xlarge"
        weighted_capacity = 4
      }

      override {
        instance_type     = "p3.16xlarge"
        weighted_capacity = 8
      }

      override {
        instance_type     = "g4dn.12xlarge"
        weighted_capacity = 4
      }
    }
  }

  # Scale based on training queue depth
  enabled_metrics = [
    "GroupDesiredCapacity",
    "GroupInServiceInstances",
    "GroupMinSize",
    "GroupMaxSize"
  ]

  tag {
    key                 = "Name"
    value               = "${var.project_name}-training-spot"
    propagate_at_launch = true
  }

  dynamic "tag" {
    for_each = local.cost_allocation_tags

    content {
      key                 = tag.key
      value               = tag.value
      propagate_at_launch = true
    }
  }
}

# Launch Template for Spot Instances
resource "aws_launch_template" "training_spot" {
  name_prefix   = "${var.project_name}-training-"
  image_id      = data.aws_ami.gpu_optimized.id
  instance_type = "p3.2xlarge"  # Default, overridden by mixed policy

  network_interfaces {
    associate_public_ip_address = false
    security_groups             = [var.training_security_group_id]
    delete_on_termination       = true
  }

  iam_instance_profile {
    arn = aws_iam_instance_profile.training_spot.arn
  }

  block_device_mappings {
    device_name = "/dev/sda1"

    ebs {
      volume_size           = 100
      volume_type           = "gp3"
      iops                  = 3000
      throughput            = 125
      delete_on_termination = true
      encrypted             = true
    }
  }

  user_data = base64encode(templatefile("${path.module}/user_data.sh", {
    cluster_name = var.eks_cluster_name
    region       = data.aws_region.current.name
  }))

  metadata_options {
    http_endpoint               = "enabled"
    http_tokens                 = "required"
    http_put_response_hop_limit = 1
  }

  tag_specifications {
    resource_type = "instance"

    tags = merge(local.cost_allocation_tags, {
      Name = "${var.project_name}-training-spot"
    })
  }

  tag_specifications {
    resource_type = "volume"

    tags = merge(local.cost_allocation_tags, {
      Name = "${var.project_name}-training-volume"
    })
  }
}

# Auto Scaling Policies
resource "aws_autoscaling_policy" "scale_up" {
  name                   = "${var.project_name}-scale-up"
  scaling_adjustment     = 2
  adjustment_type        = "ChangeInCapacity"
  cooldown               = 300
  autoscaling_group_name = aws_autoscaling_group.training_spot_fleet.name
}

resource "aws_autoscaling_policy" "scale_down" {
  name                   = "${var.project_name}-scale-down"
  scaling_adjustment     = -1
  adjustment_type        = "ChangeInCapacity"
  cooldown               = 300
  autoscaling_group_name = aws_autoscaling_group.training_spot_fleet.name
}

# CloudWatch Alarms for Auto Scaling
resource "aws_cloudwatch_metric_alarm" "training_queue_high" {
  alarm_name          = "${var.project_name}-training-queue-high"
  comparison_operator = "GreaterThanThreshold"
  evaluation_periods  = 2
  metric_name         = "ApproximateNumberOfMessagesVisible"
  namespace           = "AWS/SQS"
  period              = 60
  statistic           = "Average"
  threshold           = 5

  alarm_description = "Scale up when training queue is high"
  alarm_actions     = [aws_autoscaling_policy.scale_up.arn]

  dimensions = {
    QueueName = var.training_queue_name
  }
}

resource "aws_cloudwatch_metric_alarm" "training_queue_low" {
  alarm_name          = "${var.project_name}-training-queue-low"
  comparison_operator = "LessThanThreshold"
  evaluation_periods  = 5
  metric_name         = "ApproximateNumberOfMessagesVisible"
  namespace           = "AWS/SQS"
  period              = 60
  statistic           = "Average"
  threshold           = 2

  alarm_description = "Scale down when training queue is low"
  alarm_actions     = [aws_autoscaling_policy.scale_down.arn]

  dimensions = {
    QueueName = var.training_queue_name
  }
}

# S3 Intelligent Tiering Lifecycle Policy
resource "aws_s3_bucket_lifecycle_configuration" "intelligent_tiering" {
  bucket = var.ml_data_bucket_id

  rule {
    id     = "intelligent-tiering"
    status = "Enabled"

    filter {
      prefix = ""
    }

    transition {
      days          = 30
      storage_class = "INTELLIGENT_TIERING"
    }

    noncurrent_version_transition {
      noncurrent_days = 30
      storage_class   = "GLACIER"
    }

    noncurrent_version_expiration {
      noncurrent_days = 90
    }
  }

  rule {
    id     = "delete-old-logs"
    status = "Enabled"

    filter {
      prefix = "logs/"
    }

    expiration {
      days = 7
    }
  }

  rule {
    id     = "archive-training-data"
    status = "Enabled"

    filter {
      and {
        prefix = "training-data/"
        tags = {
          Archive = "true"
        }
      }
    }

    transition {
      days          = 90
      storage_class = "GLACIER"
    }

    transition {
      days          = 180
      storage_class = "DEEP_ARCHIVE"
    }
  }
}

# Cost Anomaly Detection
resource "aws_ce_anomaly_monitor" "ml_infrastructure" {
  name              = "${var.project_name}-cost-anomaly"
  monitor_type      = "DIMENSIONAL"
  monitor_dimension = "SERVICE"
}

resource "aws_ce_anomaly_subscription" "ml_infrastructure_alert" {
  name      = "${var.project_name}-cost-anomaly-alert"
  frequency = "DAILY"

  monitor_arn_list = [
    aws_ce_anomaly_monitor.ml_infrastructure.arn
  ]

  subscriber {
    type    = "SNS"
    address = var.cost_alert_sns_topic_arn
  }

  threshold_expression {
    and {
      dimension {
        key           = "ANOMALY_TOTAL_IMPACT_ABSOLUTE"
        match_options = ["GREATER_THAN_OR_EQUAL"]
        values        = ["100"]  # Alert on $100+ anomalies
      }
    }
  }
}

# Budget Alerts
resource "aws_budgets_budget" "ml_monthly" {
  name              = "${var.project_name}-monthly-budget"
  budget_type       = "COST"
  limit_amount      = var.monthly_budget_limit
  limit_unit        = "USD"
  time_period_start = "2024-01-01_00:00"
  time_unit         = "MONTHLY"

  cost_filter {
    name = "TagKeyValue"
    values = [
      "Project$${var.project_name}"
    ]
  }

  notification {
    comparison_operator        = "GREATER_THAN"
    threshold                  = 80
    threshold_type             = "PERCENTAGE"
    notification_type          = "ACTUAL"
    subscriber_email_addresses = var.budget_alert_emails
  }

  notification {
    comparison_operator        = "GREATER_THAN"
    threshold                  = 100
    threshold_type             = "PERCENTAGE"
    notification_type          = "ACTUAL"
    subscriber_email_addresses = var.budget_alert_emails
  }

  notification {
    comparison_operator        = "GREATER_THAN"
    threshold                  = 90
    threshold_type             = "PERCENTAGE"
    notification_type          = "FORECASTED"
    subscriber_email_addresses = var.budget_alert_emails
  }
}

# Scheduled Instance Shutdown for Development
resource "aws_lambda_function" "instance_scheduler" {
  count = var.environment != "production" ? 1 : 0

  filename      = data.archive_file.scheduler_lambda[0].output_path
  function_name = "${var.project_name}-instance-scheduler"
  role          = aws_iam_role.scheduler_lambda[0].arn
  handler       = "index.handler"
  runtime       = "python3.11"
  timeout       = 60

  environment {
    variables = {
      ENVIRONMENT  = var.environment
      PROJECT_NAME = var.project_name
    }
  }

  tags = local.cost_allocation_tags
}

# Schedule: Stop instances at night (development only)
resource "aws_cloudwatch_event_rule" "stop_instances" {
  count = var.environment != "production" ? 1 : 0

  name                = "${var.project_name}-stop-instances"
  description         = "Stop development instances at night"
  schedule_expression = "cron(0 22 * * ? *)"  # 10 PM daily
}

resource "aws_cloudwatch_event_target" "stop_instances" {
  count = var.environment != "production" ? 1 : 0

  rule      = aws_cloudwatch_event_rule.stop_instances[0].name
  target_id = "StopInstances"
  arn       = aws_lambda_function.instance_scheduler[0].arn

  input = jsonencode({
    action = "stop"
  })
}

# Schedule: Start instances in morning
resource "aws_cloudwatch_event_rule" "start_instances" {
  count = var.environment != "production" ? 1 : 0

  name                = "${var.project_name}-start-instances"
  description         = "Start development instances in morning"
  schedule_expression = "cron(0 8 * * MON-FRI *)"  # 8 AM weekdays
}

resource "aws_cloudwatch_event_target" "start_instances" {
  count = var.environment != "production" ? 1 : 0

  rule      = aws_cloudwatch_event_rule.start_instances[0].name
  target_id = "StartInstances"
  arn       = aws_lambda_function.instance_scheduler[0].arn

  input = jsonencode({
    action = "start"
  })
}

# Outputs
output "spot_fleet_asg_name" {
  value = aws_autoscaling_group.training_spot_fleet.name
}

output "cost_anomaly_monitor_arn" {
  value = aws_ce_anomaly_monitor.ml_infrastructure.arn
}

output "budget_name" {
  value = aws_budgets_budget.ml_monthly.name
}
\end{lstlisting}

\textbf{Cost Optimization Best Practices:}

\begin{itemize}
    \item \textbf{Spot Instances:} Use spot instances for 60-80\% cost savings on training workloads
    \item \textbf{Right-Sizing:} Continuously monitor and adjust instance types based on actual usage
    \item \textbf{Auto-Scaling:} Scale resources dynamically based on demand
    \item \textbf{Reserved Instances:} Purchase reserved capacity for predictable baseline workloads
    \item \textbf{Storage Tiering:} Automatically move infrequently accessed data to cheaper storage
    \item \textbf{Resource Scheduling:} Stop development resources outside business hours
    \item \textbf{Budget Alerts:} Set up proactive cost monitoring and alerts
    \item \textbf{Tagging Strategy:} Tag all resources for accurate cost allocation and tracking
\end{itemize}

\section{Practical Infrastructure Automation Exercises}

\subsection{Exercise 1: Multi-Environment ML Infrastructure}

Deploy complete ML infrastructure across development, staging, and production environments using Terraform modules. Implement:

\begin{itemize}
    \item Environment-specific resource sizing (development: minimal, production: HA)
    \item VPC with public and private subnets in multiple availability zones
    \item EKS cluster with separate node groups for training and inference
    \item S3 buckets with lifecycle policies and versioning
    \item RDS PostgreSQL with automated backups
    \item ECR repositories for container images
    \item Proper IAM roles and security groups
\end{itemize}

\textbf{Deliverables:}
\begin{itemize}
    \item Terraform module structure with separate environments
    \item Variables file for each environment
    \item State management with remote backend
    \item Infrastructure documentation
\end{itemize}

\subsection{Exercise 2: Kubernetes ML Training Helm Chart}

Create production-ready Helm chart for distributed ML training jobs. Include:

\begin{itemize}
    \item Support for PyTorch and TensorFlow distributed training
    \item Configurable GPU resource allocation
    \item Data loading from S3, GCS, and Azure Blob
    \item MLflow experiment tracking integration
    \item Health checks and pod disruption budgets
    \item Resource quotas and limits
    \item Secrets management for credentials
\end{itemize}

\textbf{Deliverables:}
\begin{itemize}
    \item Complete Helm chart with templates
    \item Values files for different model types
    \item Testing strategy for chart validation
    \item Deployment documentation
\end{itemize}

\subsection{Exercise 3: Multi-Cloud Infrastructure Deployment}

Deploy ML platform infrastructure across AWS, GCP, and Azure using Terraform. Implement:

\begin{itemize}
    \item Kubernetes clusters on EKS, GKE, and AKS
    \item Object storage on S3, GCS, and Azure Blob
    \item Managed databases for metadata storage
    \item Container registries in each cloud
    \item Cross-cloud networking for data replication
    \item Unified monitoring and logging
\end{itemize}

\textbf{Deliverables:}
\begin{itemize}
    \item Cloud-agnostic Terraform modules
    \item Provider-specific implementations
    \item Cost comparison analysis
    \item Failover and replication strategy
\end{itemize}

\subsection{Exercise 4: Automated Disaster Recovery System}

Implement automated disaster recovery for ML infrastructure. Include:

\begin{itemize}
    \item AWS Backup configuration for RDS and EBS
    \item S3 cross-region replication for models and data
    \item Automated backup verification testing
    \item Recovery time objective (RTO) and recovery point objective (RPO) monitoring
    \item Automated failover procedures
    \item Recovery testing schedule
\end{itemize}

\textbf{Deliverables:}
\begin{itemize}
    \item Terraform backup infrastructure code
    \item Python disaster recovery automation script
    \item Recovery runbook documentation
    \item Disaster recovery testing report
\end{itemize}

\subsection{Exercise 5: Cost-Optimized Training Infrastructure}

Build cost-optimized infrastructure for ML training. Implement:

\begin{itemize}
    \item Spot instance autoscaling groups with fallback to on-demand
    \item Multiple instance type support for spot diversity
    \item Training queue-based autoscaling
    \item S3 intelligent tiering for training data
    \item Scheduled shutdown of development resources
    \item Cost anomaly detection and alerting
    \item Monthly budget enforcement
\end{itemize}

\textbf{Deliverables:}
\begin{itemize}
    \item Terraform cost optimization module
    \item Cost tracking dashboard
    \item 30-day cost analysis report
    \item Optimization recommendations
\end{itemize}

\subsection{Exercise 6: Kubernetes Operator for ML Workflows}

Develop custom Kubernetes operator for ML training jobs. Create:

\begin{itemize}
    \item Custom Resource Definition (CRD) for ML training
    \item Operator controller logic in Go or Python
    \item Support for distributed training patterns
    \item Integration with MLflow for experiment tracking
    \item Automatic resource cleanup after job completion
    \item Job scheduling and prioritization
\end{itemize}

\textbf{Deliverables:}
\begin{itemize}
    \item Custom operator implementation
    \item CRD specification and examples
    \item Operator deployment manifests
    \item User guide and API documentation
\end{itemize}

\subsection{Exercise 7: Infrastructure Monitoring and Observability}

Implement comprehensive monitoring for ML infrastructure. Deploy:

\begin{itemize}
    \item Prometheus for metrics collection
    \item Grafana dashboards for infrastructure visualization
    \item Loki for log aggregation
    \item Jaeger for distributed tracing
    \item AlertManager for alert routing
    \item Infrastructure health scoring system
\end{itemize}

\textbf{Deliverables:}
\begin{itemize}
    \item Monitoring stack deployment code
    \item Custom Grafana dashboards for ML workloads
    \item Alert rules for infrastructure issues
    \item Runbook for common incidents
\end{itemize}

\section{Summary}

Infrastructure as Code is fundamental to building reliable, scalable, and cost-effective ML platforms. This chapter covered:

\textbf{1. Terraform Modules:}
\begin{itemize}
    \item Modular ML infrastructure design with reusable components
    \item EKS clusters with GPU and CPU node groups
    \item S3, RDS, and caching infrastructure
    \item Comprehensive networking and security
    \item Environment-specific configurations
\end{itemize}

\textbf{2. Kubernetes and Helm:}
\begin{itemize}
    \item Production-ready Helm charts for training and inference
    \item Distributed training with PyTorchJob
    \item Autoscaling based on workload metrics
    \item Resource quotas and pod disruption budgets
    \item Custom operators for ML workflows
\end{itemize}

\textbf{3. Multi-Cloud Management:}
\begin{itemize}
    \item AWS (EKS, S3, RDS, SageMaker)
    \item GCP (GKE, Cloud Storage, Vertex AI)
    \item Azure (AKS, Blob Storage, Azure ML)
    \item Cloud-agnostic module design
    \item Cross-cloud data replication
\end{itemize}

\textbf{4. Container Orchestration:}
\begin{itemize}
    \item Kubernetes operators for ML workloads
    \item Cluster autoscaling with priority-based expansion
    \item Resource quotas and namespace isolation
    \item GPU resource management
    \item Job scheduling and prioritization
\end{itemize}

\textbf{5. Disaster Recovery:}
\begin{itemize}
    \item Automated backup with AWS Backup
    \item Cross-region replication for critical data
    \item Recovery automation with Python
    \item RTO/RPO monitoring and enforcement
    \item Regular disaster recovery testing
\end{itemize}

\textbf{6. Cost Optimization:}
\begin{itemize}
    \item Spot instances for 60-80\% cost savings
    \item Auto-scaling based on demand
    \item Storage lifecycle management
    \item Cost anomaly detection
    \item Budget alerts and enforcement
    \item Scheduled resource management
\end{itemize}

\textbf{Key Takeaways:}

\begin{enumerate}
    \item \textbf{Automation:} Automate all infrastructure provisioning and management
    \item \textbf{Modularity:} Design reusable, composable infrastructure modules
    \item \textbf{Multi-Cloud:} Avoid vendor lock-in with cloud-agnostic designs
    \item \textbf{Cost Awareness:} Continuously monitor and optimize infrastructure costs
    \item \textbf{Resilience:} Build fault-tolerant systems with automated recovery
    \item \textbf{Security:} Implement least privilege access and encryption by default
    \item \textbf{Observability:} Monitor all infrastructure components comprehensively
\end{enumerate}

Infrastructure as Code transforms ML infrastructure from manual, error-prone processes to automated, reliable, and auditable deployments. The investment in proper IaC practices pays dividends through faster deployment cycles, reduced operational burden, and improved system reliability.
